\begin{theorem}[Jensen’s Inequality]
Let \(f:\mathbb{R}\to\mathbb{R}\) be convex, i.e.
\[
\forall\,x,y\in\mathbb{R},\;\forall\,t\in[0,1]\qquad
f\bigl(t x+(1-t)y\bigr)\le t\,f(x)+(1-t)\,f(y). \tag{1}
\]
Let \(x_{1},\dots ,x_{n}\in\mathbb{R}\) and non‑negative weights
\(\lambda_{1},\dots ,\lambda_{n}\) satisfy
\[
\lambda_{i}\ge0\;(i=1,\dots ,n),\qquad
\sum_{i=1}^{n}\lambda_{i}=1. \tag{2}
\]
Then
\[
f\!\left(\sum_{i=1}^{n}\lambda_{i}x_{i}\right)\le
\sum_{i=1}^{n}\lambda_{i}f(x_{i}). \tag{3}
\]
Equality holds iff either all the \(x_{i}\) are equal, or \(f\) is affine on
the convex hull of \(\{x_{1},\dots ,x_{n}\}\) (equivalently on the smallest
interval containing the points), or the weights attached to points where
\(f\) is not affine are zero.
\end{theorem}

\begin{proof}
We proceed by induction on the number of points \(n\).

\medskip\noindent\textbf{1. Degenerate cases.}
\begin{itemize}
\item If \(n=1\), then (2) forces \(\lambda_{1}=1\) and (3) reduces to
\(f(x_{1})\le f(x_{1})\), true with equality.
\item If some weight equals \(1\) (say \(\lambda_{k}=1\) and all other
\(\lambda_{i}=0\)), then
\(\sum_{i=1}^{n}\lambda_{i}x_{i}=x_{k}\) and
\(\sum_{i=1}^{n}\lambda_{i}f(x_{i})=f(x_{k})\); again (3) holds with
equality.
\end{itemize}
Thus the statement is valid in all degenerate situations.

\medskip\noindent\textbf{2. Base case \(n=2\).}
Write \(\lambda_{1}=t\), \(\lambda_{2}=1-t\) with \(t\in[0,1]\) (the
condition \(\lambda_{1}+\lambda_{2}=1\) guarantees this).  Applying the
definition of convexity (1) to \(x_{1},x_{2}\) yields
\[
f\bigl(t x_{1}+(1-t)x_{2}\bigr)\le t\,f(x_{1})+(1-t)f(x_{2}),
\]
which is exactly (3) for \(n=2\).

\medskip\noindent\textbf{3. Induction hypothesis.}
Assume that for some integer \(k\ge 2\) the inequality (3) holds for
\emph{every} collection of at most \(k\) points with non‑negative weights
summing to \(1\).  We shall prove it for \(n=k+1\).

\medskip\noindent\textbf{4. Inductive step.}
Let \(x_{1},\dots ,x_{n}\) (\(n=k+1\)) and weights \(\lambda_{1},\dots ,\lambda_{n}\)
satisfy (2).  Set
\[
L:=\sum_{i=1}^{n-1}\lambda_{i}. \tag{4}
\]
If \(L=0\) then \(\lambda_{n}=1\) and we are back to a degenerate case.
Hence assume \(L>0\).

Define normalized weights for the first \(n-1\) indices:
\[
\mu_{i}:=\frac{\lambda_{i}}{L},\qquad i=1,\dots ,n-1. \tag{5}
\]
Clearly \(\mu_{i}\ge0\) and \(\sum_{i=1}^{n-1}\mu_{i}=1\).

Set
\[
y:=\sum_{i=1}^{n-1}\mu_{i}x_{i}. \tag{6}
\]
By the induction hypothesis applied to the \((n-1)\)-tuple
\((x_{1},\dots ,x_{n-1})\) with weights \(\mu_{i}\) we have
\[
f(y)\le\sum_{i=1}^{n-1}\mu_{i}f(x_{i}). \tag{7}
\]

The original convex combination can be rewritten as a two‑point
combination of \(y\) and \(x_{n}\):
\begin{align*}
\sum_{i=1}^{n}\lambda_{i}x_{i}
   &=\sum_{i=1}^{n-1}\lambda_{i}x_{i}+\lambda_{n}x_{n} \\
   &=L\Bigl(\sum_{i=1}^{n-1}\mu_{i}x_{i}\Bigr)+\lambda_{n}x_{n}
    =L\,y+\lambda_{n}x_{n}. \tag{8}
\end{align*}
From (2) and (4) we obtain
\[
L+\lambda_{n}=1. \tag{9}
\]

Apply the two‑point convexity (1) to the points \(y\) and \(x_{n}\) with
weight \(L\):
\[
f(Ly+\lambda_{n}x_{n})\le L\,f(y)+\lambda_{n}f(x_{n}). \tag{10}
\]
Substituting (8) for the argument of \(f\) and using (7) to bound \(f(y)\)
gives
\begin{align*}
f\!\left(\sum_{i=1}^{n}\lambda_{i}x_{i}\right)
   &\le L\!\left(\sum_{i=1}^{n-1}\mu_{i}f(x_{i})\right)+\lambda_{n}f(x_{n}) \\
   &=\sum_{i=1}^{n-1}L\mu_{i}f(x_{i})+\lambda_{n}f(x_{n}) \\
   &=\sum_{i=1}^{n-1}\lambda_{i}f(x_{i})+\lambda_{n}f(x_{n}),
\end{align*}
which is precisely (3) for \(n\) points.  Hence the inequality holds for
\(n=k+1\).  By induction it holds for every integer \(n\ge1\).

\medskip\noindent\textbf{5. Equality case.}
Equality in (1) occurs exactly when either the weight \(t\) equals \(0\) or
\(1\) (so the convex combination reduces to a single point) or the
restriction of \(f\) to the segment \([x,y]\) is affine:
\[
f\bigl(t x+(1-t)y\bigr)=t\,f(x)+(1-t)f(y)\quad\forall\,t\in[0,1].
\]

In the inductive proof equality in the final inequality requires equality
at the two points where we used (1) (step (10)) and where we invoked the
induction hypothesis (step (7)).  Consequently equality in (3) holds
iff one of the following equivalent conditions is satisfied:

\begin{enumerate}
\item All points coincide: \(x_{1}=x_{2}=\dots =x_{n}\).  Then both sides of
(3) equal \(f(x_{1})\).
\item \(f\) is affine on the convex hull of the points, i.e. there exist
constants \(a,b\) such that \(f(t)=a t+b\) for every \(t\) in the smallest
interval containing \(\{x_{1},\dots ,x_{n}\}\).  On such an interval the
inequality (1) is an equality for any weight, and every induction step
becomes an equality.
\item The weights attached to points where \(f\) is not affine are zero.
Equivalently, all points with positive weight lie in an interval on which
\(f\) is affine.
\end{enumerate}
These are precisely the conditions stated in the theorem. \(\square\)
\end{proof}